% Clue Selection From a Learned Listener
%
% Build:  tectonic docs/clue-selection-learned.tex
% Output: docs/clue-selection-learned.pdf
%
% Companion to clue-selection-theory.tex. Same problem, same rewards, same
% notation; the guesser model is a learned scoring function rather than
% Gaussian noise on similarities, and the resulting expectations turn out to
% be cheaper rather than dearer.

\documentclass[12pt]{article}

\usepackage[a4paper,margin=3.2cm]{geometry}
\usepackage{amsmath}
\usepackage{amssymb}
\usepackage{amsthm}
\usepackage{booktabs}
\usepackage{caption}
\usepackage{float}
\usepackage[skip=8pt, indent=0pt]{parskip}
\usepackage{microtype}
\usepackage[T1]{fontenc}
\usepackage{libertinus}

\DeclareMathOperator*{\argmin}{arg\,min}

\newtheorem{lemma}{Lemma}

\begin{document}

\begin{center}
  {\LARGE Clue Selection From a Learned Listener}\\[0.7em]
  Codenames spymaster, exact expected reward under a learned choice model
\end{center}
\vspace{1.5em}

\section{Setup}
\label{sec:setup}

As before, let $A$ be the set of our team's unrevealed board words and $B$ the
set of all other unrevealed board words, with $i$ indexing $A$ and $w$
indexing $B$. Selecting a word of $A$ is worth $1$, and selecting $w \in B$
costs $c_w > 0$. We choose a clue word and a number $k$, with $1 \le k \le K$,
to maximize expected reward.

In place of a noise model on similarities, suppose we have a model that, for a
fixed clue word, assigns each board word $x$ a real score $s_x$, and predicts
that a guesser choosing from a set $S$ of unrevealed words picks $x \in S$ with
probability
\begin{equation}
  \label{eq:pl}
  \Pr(\text{pick } x \mid S) \;=\; \frac{\lambda_x}{\sum_{y \in S} \lambda_y},
  \qquad \lambda_x \;=\; e^{s_x}.
\end{equation}
The guesser selects words one at a time under \eqref{eq:pl}, removing each
selection from $S$, and stops after $k$ selections or upon selecting a word of
$B$, whichever comes first.

Two things are worth making explicit about \eqref{eq:pl}. First, the scores do
not depend on $S$: one score per word is computed once, and removing a word
changes only the normalizing sum. This is Luce's choice axiom, and the model
that supplies the scores is fitted under exactly this assumption, so using it
here is a matter of consistency rather than of convenience. Second,
\eqref{eq:pl} applied repeatedly defines a distribution over orderings known as
the Plackett--Luce model, and the expected reward we want is a sum over that
distribution: each first selection opens a subtree of possible second
selections with renormalized probabilities, and so on. Enumerating that tree
costs $O(|A \cup B|^{K})$. Section~\ref{sec:race} removes the tree entirely.

\section{The exponential race}
\label{sec:race}

\begin{lemma}
  \label{lem:race}
  Let $T_x \sim \mathrm{Exp}(\lambda_x)$ be independent, one per word. Then the
  ordering of the words by increasing $T_x$ has exactly the distribution
  generated by repeated selection under \eqref{eq:pl}.
\end{lemma}

\begin{proof}
  For independent exponentials, $\Pr\bigl(\argmin_x T_x = x_1\bigr) =
  \lambda_{x_1} / \sum_y \lambda_y$, which is \eqref{eq:pl} for the first
  selection. By the memorylessness of the exponential, conditional on
  $T_{x_1} = t$ the remaining variables $T_x - t$ are again independent
  exponentials with the same rates $\lambda_x$, so the same statement applies
  to the second selection with $x_1$ removed, and inductively to the rest.
  Multiplying the successive factors,
  \[ \Pr\bigl(T_{x_1} < T_{x_2} < \cdots < T_{x_n}\bigr)
     \;=\; \prod_{j=1}^{n} \frac{\lambda_{x_j}}{\sum_{l \ge j} \lambda_{x_l}}, \]
  which is the probability the repeated-selection process produces that
  ordering. \qedhere
\end{proof}

Lemma~\ref{lem:race} is an identity, not an approximation: the race reproduces
the selection tree at every node. What it buys is that the quantities we need
are now functions of minima of independent exponentials, which are elementary.

Following the same shape as the noise model, let
\[ T \;=\; \min_{w \in B} T_w, \qquad
   W \;=\; \argmin_{w \in B} T_w, \qquad
   N \;=\; \#\{\, i \in A : T_i < T \,\}. \]
So $W$ is the word of $B$ the guesser would reach first, $T$ is when, and $N$
is how many of our words they would take before it. The guesser reveals
$\min(k, N)$ words of $A$, and selects a word of $B$ --- necessarily $W$ ---
if and only if $N < k$. The objective is to maximize $G(k) - P(k)$ for
\[ G(k) = \mathbb E\bigl[\min(k, N)\bigr], \qquad
   P(k) = \mathbb E\bigl[\, c_W \, \mathbf 1\{N < k\} \,\bigr]. \]

\section{Derivation}
\label{sec:derivation}

Write $\Lambda = \sum_{w \in B} \lambda_w$.

\paragraph{The distribution of $T$.} A minimum of independent exponentials is
exponential with the summed rate, so $T \sim \mathrm{Exp}(\Lambda)$ and
\begin{equation}
  \label{eq:T}
  F_T(t) = 1 - e^{-\Lambda t}, \qquad dF_T(t) = \Lambda e^{-\Lambda t}\, dt .
\end{equation}

\paragraph{The identity of $W$.} For independent exponentials the minimum and
the argument of the minimum are independent, and
\begin{equation}
  \label{eq:W}
  \Pr(W = w) \;=\; \frac{\lambda_w}{\Lambda}
  \qquad\text{for every } t .
\end{equation}
This is where the learned model is cheaper than the noise model. There, the
identity of the first non-team word depended on $t$ and required a hazard
ratio at every quadrature point; here it does not depend on $t$ at all, so
\[ \bar c \;=\; \mathbb E[c_W] \;=\; \frac{1}{\Lambda}\sum_{w \in B} c_w \lambda_w \]
is a single constant, computed once per clue in $O(|B|)$.

\paragraph{The distribution of $N$ given $T$.} $T$ is a function of
$(T_w)_{w \in B}$ only, which is independent of $(T_i)_{i \in A}$. Conditioning
on $T = t$ therefore leaves each $T_i$ exponential with rate $\lambda_i$, and
each $i \in A$ independently satisfies $T_i < t$ with probability
\begin{equation}
  \label{eq:p}
  p_i(t) \;=\; 1 - e^{-\lambda_i t}.
\end{equation}
So given $T = t$, $N$ is a sum of $|A|$ independent Bernoulli variables with
success probabilities $p_i(t)$ --- a Poisson binomial, as in the noise model.

\paragraph{Assembling.} Using $\min(k,N) = \sum_{j=1}^{k}\mathbf 1\{N \ge j\}$
and conditioning on $T$,
\begin{equation}
  \label{eq:G}
  G(k) \;=\; \sum_{j=1}^{k} \int_0^\infty \Pr(N \ge j \mid T = t)\, dF_T(t).
\end{equation}
For $P(k)$, note that given $T = t$ the count $N$ depends only on
$(T_i)_{i \in A}$ while $W$ depends only on $(T_w)_{w \in B}$, so $c_W$ and
$\mathbf 1\{N < k\}$ are conditionally independent given $T$. Combined with
\eqref{eq:W}, which makes $\mathbb E[c_W \mid T = t]$ constant in $t$, the
expectation factors completely:
\begin{equation}
  \label{eq:P}
  P(k) \;=\; \bar c \int_0^\infty \Pr(N < k \mid T = t)\, dF_T(t).
\end{equation}

\section{Numerical evaluation}
\label{sec:computation}

Both \eqref{eq:G} and \eqref{eq:P} have the form $\int_0^\infty \psi(t)\,
dF_T(t)$ for a bounded $\psi$. Unlike the noise model, no truncation is needed:
substituting $u = e^{-\Lambda t}$ gives $du = -\Lambda e^{-\Lambda t} dt$, so
$dF_T(t) = -du$, and as $t$ runs over $(0, \infty)$, $u$ runs over $(0,1)$.
Hence
\begin{equation}
  \label{eq:sub}
  \int_0^\infty \psi(t)\, dF_T(t) \;=\; \int_0^1 \psi\bigl(t(u)\bigr)\, du,
  \qquad t(u) = -\frac{\ln u}{\Lambda},
\end{equation}
an integral of a bounded function over the unit interval, which a midpoint rule
on a uniform grid handles without any choice of bounds. The whole tail of $T$
is carried exactly by the substitution rather than being cut off and argued to
be negligible.

The integrand is bounded but not smooth at $u = 0$: by \eqref{eq:pu} it carries
$u^{\lambda_i/\Lambda}$, whose derivative is unbounded there whenever that
ratio is below $1$. The midpoint rule therefore converges at roughly first
order rather than second, and the remedy is simply more cells, since each one
costs a power and nothing else. In practice far fewer are needed than the
error alone suggests, because the search consumes the ordering of clues rather
than their values, and the ordering stabilizes well before the values do.

The success probabilities take a convenient form in $u$: from \eqref{eq:p},
\begin{equation}
  \label{eq:pu}
  p_i(u) \;=\; 1 - u^{\lambda_i / \Lambda},
\end{equation}
so each quadrature point costs one power per word of $A$ and no special
functions anywhere in the computation.

The conditional probabilities $\Pr(N \ge j \mid T = t)$ follow from the
standard recursion for a sum of independent Bernoulli variables. Letting
$q^{(i)}_r$ be $\Pr(N = r)$ after the first $i$ words of $A$,
\[ q^{(0)}_0 = 1, \qquad
   q^{(i)}_r \;=\; q^{(i-1)}_r\bigl(1 - p_i\bigr) \;+\; q^{(i-1)}_{r-1}\,p_i , \]
and $\Pr(N \ge j) = 1 - \sum_{r < j} q_r$. Counts above $K$ are never used, so
the recursion is truncated there and costs $O(|A|K)$ per quadrature point.

With $M$ cells the total is $O\bigl(M|A|K + |B|\bigr)$ per clue, against
$O\bigl(M(|A|K + |B|)\bigr)$ for the noise model: the $|B|$ term leaves the
integral because of \eqref{eq:W}. Every value of $k$ from $1$ to $K$ is
obtained from one pass, since \eqref{eq:G} accumulates over $j$ and
\eqref{eq:P} needs the same truncated distribution of $N$.

\section{What is exact and what is assumed}
\label{sec:assumptions}

It is worth separating the two, because the derivation above is exact in a way
that can disguise where the modelling risk actually sits.

Exact, given \eqref{eq:pl}: Lemma~\ref{lem:race}, so the race is the selection
tree rather than a stand-in for it; the distribution \eqref{eq:T} of $T$; the
independence and law \eqref{eq:W} of $W$; the conditional Poisson binomial
\eqref{eq:p}; the factorization of $P(k)$; the substitution \eqref{eq:sub}; and
the Bernoulli recursion. The cap $\min(k, N)$ is likewise exact and not a
truncation, since $N \ge k$ is precisely the event that the guesser exhausts
its guesses on our words and never reaches $W$. The sole numerical error is the
quadrature in \eqref{eq:sub}, over a bounded integrand on a bounded interval.

Assumed, and not established here:

\emph{Scores do not depend on the remaining set.} This is the content of
\eqref{eq:pl} and the one substantive modelling commitment. It is not a cost of
the race construction --- a direct enumeration of the selection tree assumes
the same thing the moment it renormalizes the scores over the survivors rather
than recomputing them. If the guesser's scores genuinely change as the board
empties, both routes are wrong together, and Plackett--Luce is the wrong family
rather than the integral being the wrong evaluation of it.

\emph{The guesser spends every guess.} The process stops only on reaching $k$
selections or a word of $B$. A guesser who declines a guess it is unsure of
faces a different expected reward, and nothing in the listener model expresses
that choice.

\emph{The clue number is the guess budget.} $k$ enters only as the cap. Where
the rules grant a bonus guess, $K$ should be set to the true budget rather
than to the announced number.

\end{document}
